\documentclass[11pt]{article}
\usepackage[letterpaper,margin=0.68in,headheight=15pt]{geometry}
\usepackage[T1]{fontenc}
\usepackage{lmodern}
\usepackage{amsmath,amssymb}
\usepackage{array,booktabs,enumitem,multicol}
\usepackage{xcolor}
\usepackage{tikz}
\usepackage{fancyhdr}
\usepackage{microtype}
\setlength{\parindent}{0pt}
\setlength{\parskip}{5pt}
\definecolor{olive}{HTML}{555B35}
\definecolor{gold}{HTML}{B18A22}
\definecolor{soft}{HTML}{F4F1DE}
\definecolor{bluegray}{HTML}{E8EEF2}
\pagestyle{fancy}
\fancyhf{}
\fancyhead[L]{\textcolor{olive}{\small MATH\& 141: Precalculus I}}
\fancyhead[R]{\textcolor{gold}{\small Fall 2026}}
\fancyfoot[C]{\thepage}
\renewcommand{\headrulewidth}{0.45pt}
\renewcommand{\footrulewidth}{0pt}
\setlist[enumerate]{leftmargin=*,itemsep=0.08in,topsep=0.05in}
\newcommand{\assessmenttitle}[2]{%
  {\Large\bfseries Module #1: #2}\par\vspace{3pt}
}
\newcommand{\studentline}{%
\noindent\textbf{Name:}\hspace{2.6in}\textbf{Date:}\hspace{1.15in}\par\vspace{7pt}}
\newcommand{\directions}[1]{\textit{#1}\par\vspace{4pt}}
\newcommand{\answerblank}[1]{\vspace{#1}}
\newcommand{\blankgrid}[4]{%
\begin{center}
\begin{tikzpicture}[x=0.75cm,y=0.75cm]
  \draw[step=1,gray!28,very thin] (#1,#3) grid (#2,#4);
  \draw[->,thick] (#1,0)--(#2,0) node[right] {$x$};
  \draw[->,thick] (0,#3)--(0,#4) node[above] {$y$};
  \foreach \x in {#1,...,#2}{\ifnum\x=0\else\draw (\x,0.10)--(\x,-0.10) node[below=2pt,font=\scriptsize]{\x};\fi}
  \foreach \y in {#3,...,#4}{\ifnum\y=0\else\draw (0.10,\y)--(-0.10,\y) node[left=2pt,font=\scriptsize]{\y};\fi}
\end{tikzpicture}
\end{center}}
\newcommand{\smallgrid}[4]{%
\begin{tikzpicture}[x=0.50cm,y=0.50cm]
  \draw[step=1,gray!28,very thin] (#1,#3) grid (#2,#4);
  \draw[->,thick] (#1,0)--(#2,0) node[right] {$x$};
  \draw[->,thick] (0,#3)--(0,#4) node[above] {$y$};
\end{tikzpicture}}
\begin{document}

% MODULE 3 PREQUIZ
\assessmenttitle{3}{Polynomials and Their Shape -- Prequiz}
\studentline
\directions{Four short questions. Use degree, leading term, roots, and multiplicity as appropriate.}
\begin{enumerate}
\item Determine the degree and end behavior of
\[
h(x)=-3x^6+2x^3-5.
\]
\answerblank{0.55in}
\item For
\[
p(x)=(x+2)^3(x-1)^2,
\]
list the real roots and their multiplicities. State whether the graph crosses or touches the $x$-axis at each root.
\answerblank{0.65in}
\item What is the maximum possible number of turning points of a degree-$5$ polynomial?
\answerblank{0.38in}
\item Describe the end behavior of a polynomial with positive leading coefficient and odd degree.
\answerblank{0.50in}
\end{enumerate}
\newpage

% MODULE 3 CHECKIN PAGE 1
\assessmenttitle{3}{Polynomials and Their Shape -- Weekly Check-In}
\studentline
\directions{Three questions. Use degree, leading term, roots, and multiplicity as appropriate.}
\textbf{1. Degree and end behavior.}

For
\[
f(x)=-2x^5+3x^2-7,
\]
state the degree, identify the leading term, and describe the end behavior as $x\to-\infty$ and $x\to\infty$.
\answerblank{1.15in}

\textbf{2. Roots and multiplicity.}

For
\[
p(x)=(x+2)^2(x-1)^3(x-4),
\]
list each real root and its multiplicity. At each root, state whether the graph crosses or touches the $x$-axis. What is the maximum possible number of turning points of $p$?
\answerblank{1.65in}
\newpage

% MODULE 3 CHECKIN PAGE 2
\assessmenttitle{3}{Polynomials and Their Shape -- Weekly Check-In, continued}
\studentline
\textbf{3. Sketching from factored form.}

Consider
\[
h(x)=-(x+3)(x+1)^2(x-2)^3.
\]
Before graphing, state the degree and end behavior. Then sketch a possible graph of $h(x)$ below, showing the correct behavior at every real zero. Exact vertical scale is not required.
\answerblank{0.55in}
\blankgrid{-5}{4}{-5}{5}

\end{document}
